\documentclass[12pt]{article}

\usepackage{jheppub}
\usepackage{iip-arxiv}
\toccontinuoustrue % Não quebra a página título
\compress % Comprime o resumo e sumário para caber em uma página

\usepackage{amsmath}
\usepackage{amssymb}
\usepackage{braket}
\usepackage{xcolor}

% \title{Fourier Analysis of VQE in Quantum Chemistry: An Experimental Approach}
\title{On the Fourier Structure of Energy Landscapes in Variational Quantum Algorithms for Molecular Systems}

\author[a]{Yann K. J. Leão}
\affiliation[a,b]{Federal Rural University of Pernambuco}
\emailAdd{yann.kjleao@ufrpe.br}
\author[b]{and Fábio Novaes}
\emailAdd{fabio.novaes@ufrpe.br}


\date{\today}
\abstract{Variational Quantum Algorithms (VQAs) are widely used for estimating ground-state energies of molecular systems, but their optimization remains challenging due to the non-convex structure of the energy landscape. In this work, we investigate the Fourier structure of VQE energy landscapes for molecular systems and uncover a simple but robust pattern: along one-dimensional directions in parameter space, the energy is consistently dominated by a single harmonic component near the optimum. We demonstrate this behavior empirically for H$_2$, LiH, and BeH$_2$, using parity mapping with $\mathbb{Z}_2$ tapering, where the Fourier spectrum exhibits near-zero spectral entropy, indicating an effectively monofrequency structure. Leveraging this observation, we propose a Fourier-informed line-search strategy that analytically estimates the minimum along a given direction using first-order Fourier coefficients. Our results show that this approach improves the accuracy of VQE optimization compared to standard methods. These findings suggest that, despite the apparent complexity of VQE landscapes, their local structure is highly regular and can be exploited to design more efficient optimization strategies, opening new opportunities for integrating structured models into quantum chemistry simulations.}

\begin{document}
\maketitle

% ============================================================
\section{Introduction}

% TODO:
% - Contexto: VQE para química
% - Problema: custo de otimização e medições
% - Observação: estrutura do landscape pouco explorada
% - Ideia do paper: analisar via Fourier
% - Contribuições (bullet points curtos)

% ============================================================

% Quantum chemistry is a field that has seen significant advancements with the development of variational quantum algorithms, particularly the Variational Quantum Eigensolver (VQE). The VQE algorithm is designed to find the ground state energy of a molecular system by minimizing an objective function that depends on the parameters of a quantum circuit. This approach combines classical optimization techniques with quantum computing capabilities to solve complex problems in chemistry. However, the performance of VQE can be influenced by various factors, including the choice of ansatz, the optimization landscape, and the presence of noise in quantum hardware. To better understand these factors, we can apply Fourier analysis to the VQE objective function, which allows us to decompose it into its frequency components. This decomposition can provide insights into the structure of the optimization landscape and help identify potential challenges in finding the global minimum. In this experimental approach, we will explore the Fourier analysis of VQE in the context of quantum chemistry, using numerical simulations to illustrate the concepts and techniques involved.

% Fourier analysis has been tied to Quantum Machine Learning (QML) in the past, as in~\cite{schuldEffectDataEncoding2021}, where the authors show that quantum models can be expressed as Fourier series. In this work, we will apply similar techniques to analyze the VQE objective function, which is a crucial component of the algorithm. By understanding the frequency components of the objective function, we can gain insights into the optimization process and potentially improve the performance of VQE in quantum chemistry applications.

% \section{Objectives}

% The main objectives of this study are:
% \begin{itemize}
%     \item To apply Fourier analysis to the VQE objective function in quantum chemistry.
%     \item To investigate the impact of different ansatz choices on the frequency components of the objective function.
%     \item To analyze the optimization landscape of VQE and identify potential challenges in finding the global minimum.
%     \item To explore the effects of noise on the Fourier components and the optimization process.
% \end{itemize}

Variational Quantum Algorithms (VQAs) have emerged as a leading approach for estimating ground-state energies of molecular systems on near-term quantum devices. Despite their success in small-scale applications, their optimization remains challenging due to the non-convex structure of the underlying energy landscape.

Recent theoretical works have analyzed the convergence properties of VQEs from a geometric perspective, showing that favorable optimization behavior can arise under certain conditions on the ansatz structure. In parallel, it is known that parameterized quantum circuits admit Fourier representations with respect to their parameters, suggesting that the energy landscape may possess exploitable spectral structure.

In this work, we investigate the Fourier structure of VQE energy landscapes for molecular systems. Our key observation is that, along one-dimensional cuts in parameter space, the energy is dominated by a single Fourier component in the vicinity of the optimum. This behavior is consistently observed across different molecules, including H$_2$, LiH, and BeH$_2$.

Based on this observation, we propose a simple Fourier-informed optimization step, where the minimum along a given direction is estimated analytically from the first harmonic. We demonstrate that this approach improves the accuracy of VQE results.

Our contributions are:
\begin{itemize}
    \item We empirically characterize the Fourier structure of VQE energy landscapes for molecular systems.
    \item We show that local energy variations are well approximated by a single harmonic component.
    \item We propose a Fourier-based line-search method that improves VQE optimization.
\end{itemize}

\section{Background}

\subsection{Variational Quantum Algorithms}

\begin{equation}
E(\theta) = \langle \psi(\theta) | H | \psi(\theta) \rangle
\end{equation}

% TODO:
% - explicar VQE rapidamente
% - mencionar Hamiltoniano em Pauli (sem detalhes excessivos)

\section{Local Harmonic Structure of VQE Landscapes}

Consider a parameterized quantum circuit of the form
\begin{equation}
U(\theta) = e^{-i \theta G},
\end{equation}
where $G$ is a Hermitian generator (typically a Pauli operator or a linear combination thereof).

The energy expectation value is given by
\begin{equation}
E(\theta) = \langle \psi | U^\dagger(\theta) H U(\theta) | \psi \rangle.
\end{equation}

Using standard operator identities, this expression can be written as
\begin{equation}
E(\theta) = a + b \cos\theta + c \sin\theta,
\end{equation}
for some real coefficients $a$, $b$, and $c$.

Therefore, along any single-parameter direction, the energy is naturally expressed as a first-order trigonometric polynomial. In multi-parameter settings, higher-order Fourier components may appear globally, but in the vicinity of a local optimum these higher-order terms are typically suppressed due to the smoothness of the energy functional.

As a result, the local energy landscape can be approximated as
\begin{equation}
E(\theta) \approx E_0 + A \cos(\theta - \phi),
\end{equation}
which corresponds to a single dominant Fourier mode.


\subsection{Fourier Representation}

\begin{equation}
E(\theta) = \sum_k c_k e^{ik\theta}
\end{equation}

% TODO:
% - ideia intuitiva: função periódica
% - não entrar em teoria profunda

% ============================================================
\section{Analytical Example: H$_2$}

% TODO:
% - escrever forma do Hamiltoniano (breve)
% - mostrar resultado final:

\begin{equation}
E(\theta) = a + b\cos(\theta) + c\sin(\theta)
\end{equation}

% TODO:
% - explicar que poucos termos aparecem
% - destacar sparsidade espectral

% ============================================================
\section{Methodology}

% TODO:
% - PySCF para Hamiltoniano
% - Qiskit para VQE
% - Ansatz utilizado (UCC ou EfficientSU2)

We study the Fourier structure of VQE energy landscapes for the molecules H$_2$, LiH, and BeH$_2$, using parity mapping with $\mathbb{Z}_2$ tapering.

For each system, we perform the following procedure:

\begin{enumerate}
    \item Run VQE to obtain an approximate optimal parameter vector $\mathbf{p}^*$.
    \item Select a direction $\mathbf{v}$ in parameter space.
    \item Define a one-dimensional cut:
    \begin{equation}
    \mathbf{p}(\theta) = \mathbf{p}^* + \theta \mathbf{v}.
    \end{equation}
    \item Evaluate the energy $E(\theta)$ over $\theta \in [0, 2\pi)$.
    \item Fit a first-order Fourier model:
    \begin{equation}
    E(\theta) \approx \frac{a_0}{2} + a_1 \cos\theta + b_1 \sin\theta.
    \end{equation}
\end{enumerate}

The predicted minimum along the direction is then obtained analytically as
\begin{equation}
\theta^* = \mathrm{atan2}(b_1, a_1) + \pi.
\end{equation}

This procedure defines a Fourier-informed line-search step.

\subsection{Parameter Scan}

% TODO:
% - fixar todos parâmetros menos 1
% - varrer θ em [0, 2π]

\subsection{Fourier Extraction}

% TODO:
% - FFT ou regressão
% - obter coeficientes c_k

\subsection{Spectral Metric}

\begin{equation}
p_k = \frac{|c_k|^2}{\sum_j |c_j|^2}
\end{equation}

\begin{equation}
H = -\sum_k p_k \log p_k
\end{equation}

% TODO:
% - dizer que é uma métrica simples de dispersão

% ============================================================
\section{Results}

We applied the proposed methodology to H$_2$, LiH, and BeH$_2$. In all cases, the Fourier analysis revealed that the energy landscape along sampled directions is dominated by a single harmonic component.

This is reflected in the spectral entropy, which was consistently close to zero across all systems, indicating that a single frequency captures almost all the variation in energy.

\subsection{Local Fourier Structure}

Figure~X shows the Fourier spectrum of the energy along representative directions. In all cases, the first harmonic ($k=1$) dominates, with negligible contribution from higher-order components.

\subsection{Fourier-based Optimization}

Using the fitted Fourier coefficients, we computed the predicted minimum along each direction. This prediction was found to be highly accurate, leading to improved energy estimates when used as an update step in the VQE procedure.

Compared to standard optimization, the Fourier-informed approach achieved lower final errors and improved convergence behavior.

\subsection{Discussion}

The observed monofrequency structure suggests that the local VQE landscape is effectively harmonic, even for systems with many Hamiltonian terms such as LiH and BeH$_2$.

This provides an explanation for the favorable optimization behavior observed in practice and indicates that analytical structure can be exploited to design more efficient optimization strategies.

\subsection{H$_2$}

% TODO:
% - plot E(θ)
% - plot reconstrução Fourier
% - comentar poucos coeficientes

\subsection{LiH}

% TODO:
% - mesmo procedimento
% - destacar diferença de escala

\subsection{Spectral Analysis}

% TODO:
% - gráfico dos coeficientes
% - mostrar dominância de poucos modos

% \begin{figure}
% \centering
% \includegraphics[width=0.6\textwidth]{spectrum.png}
% \caption{Fourier spectrum of the energy landscape.}
% \end{figure}

\section{Related Work}

Recent work has analyzed the convergence of VQE from a geometric perspective, showing that favorable optimization behavior can arise under conditions such as local surjectivity of the ansatz. While these results provide theoretical guarantees, they do not explicitly characterize the functional form of the energy landscape.

On the other hand, it is known that parameterized quantum circuits admit Fourier representations, where the frequency components are determined by the generator structure of the circuit. However, these works do not investigate the empirical structure of VQE energy landscapes in molecular systems.

In contrast, our work provides an empirical characterization of the Fourier structure of VQE energy landscapes and demonstrates that, in practice, they are dominated by a single harmonic component. We further show how this structure can be exploited to improve optimization.

% ============================================================
\section{Discussion}

% TODO:
% - interpretar sparsity
% - ligação com compressão de funções
% - conexão leve com IA:
%   - aprendizado
%   - modelos reduzidos

% TODO:
% - evitar teoria pesada aqui

% ============================================================
\section{Conclusion}

% TODO:
% - resumo dos achados
% - frase forte:
%   energy landscapes são dominados por baixas frequências

% TODO:
% - trabalho futuro:
%   - mais parâmetros
%   - outros sistemas
%   - integração com ML

% ============================================================


\bibliographystyle{JHEP}
\bibliography{qth-chemistry-vqe}

\end{document}
